\chapter{Everyday Commands}
\label{ch:commands}

\chapabstract{The commands you will type dozens of times a day, grouped
by task, with the flags that actually matter. Everything here is meant
to be looked up, not memorised in one sitting.}

\section{One-time setup}

Run once per machine, so every commit is correctly attributed.

\begin{shcode}
git config --global user.name  "Luigi Giuffrida"
git config --global user.email "luigi.giuffrida@polito.it"
git config --global init.defaultBranch main
git config --global core.editor "nano"   # or vim, code --wait, ...
\end{shcode}

\section{Starting a repository}

\begin{tabularx}{\linewidth}{@{}l>{\raggedright\arraybackslash}X@{}}
\toprule
\textbf{Command} & \textbf{What it does}\\
\midrule
\code{git init} & Turns the current folder into a new, empty repository.\\
\code{git clone <url>} & Copies an existing remote repository, including
  its full history, into a new local folder.\\
\code{git clone <url> <dir>} & Same, but names the local folder
  \code{<dir>} instead of using the remote's default name.\\
\bottomrule
\end{tabularx}

\section{The basic loop: status, add, commit}

\begin{shcode}
git status                 # what changed, what is staged
git diff                   # unstaged changes, line by line
git diff --staged          # staged changes, line by line
git add file1.sv file2.sv  # stage specific files
git add .                  # stage everything in and below the current dir
git commit -m "Fix reset polarity in cordic_rot"
git commit -am "message"   # stage all tracked, modified files, then commit
\end{shcode}

\begin{tipbox}{Write commit messages that will help future-you}
Use the imperative mood ("Fix", not "Fixed" or "Fixes") and keep the
first line under about 50 characters. If you need to explain why, leave
a blank line and write a longer body below. \code{git log} then shows a
clean one-line summary per commit.
\end{tipbox}

\begin{gotcha}{\code{git add .} adds everything, including junk}
Untracked build artefacts, editor swap files and simulator logs get
staged along with real changes. Add a \file{.gitignore} before your
first commit (\cref{sec:commands:gitignore}), and check \code{git
status} before every commit.
\end{gotcha}

\section{Inspecting history}

\begin{tabularx}{\linewidth}{@{}l>{\raggedright\arraybackslash}X@{}}
\toprule
\textbf{Command} & \textbf{What it does}\\
\midrule
\code{git log} & Full commit history, newest first.\\
\code{git log --oneline --graph --all} & Compact view with the branch
  topology; the single most useful \code{log} incantation.\\
\code{git show <commit>} & The full diff introduced by one commit.\\
\code{git blame <file>} & Which commit last touched each line.\\
\code{git diff <a>..<b>} & Everything that changed between two commits
  or branches.\\
\bottomrule
\end{tabularx}

\section{Branching and merging}

\begin{shcode}
git branch                    # list local branches
git branch feature/uart       # create a branch (does not switch to it)
git switch feature/uart       # switch to it
git switch -c feature/uart    # create and switch in one step
git switch main
git merge feature/uart        # bring feature/uart's commits into main
git branch -d feature/uart    # delete, once merged
\end{shcode}

\begin{notebox}{\code{switch}/\code{restore} vs. \code{checkout}}
\code{git checkout} used to do both "switch branch" and "discard a
file's changes", which was a frequent source of accidents. Modern
\Git{} (2.23+) splits these into \code{git switch} (branches) and
\code{git restore} (files). Both forms work; prefer the new ones.
\end{notebox}

\subsection*{Merge conflicts}

A conflict happens when the same lines were changed differently on both
branches. \Git{} stops the merge and marks the file:

\begin{txtcode}
<<<<<<< HEAD
q <= d;
=======
q <= d_reg;
>>>>>>> feature/uart
\end{txtcode}

Edit the file to the version you want, remove the marker lines, then:

\begin{shcode}
git add conflicted_file.sv
git commit          # completes the merge
\end{shcode}

\section{Undoing things}

Pick the row that matches how far the change has gone.

\begin{tabularx}{\linewidth}{@{}l>{\raggedright\arraybackslash}X@{}}
\toprule
\textbf{Situation} & \textbf{Command}\\
\midrule
Discard unstaged edits to a file & \code{git restore <file>}\\
Unstage a file, keep the edits & \code{git restore --staged <file>}\\
Fix the message of the last commit & \code{git commit --amend}\\
Undo the last commit, keep the changes staged &
  \code{git reset --soft HEAD\textasciitilde 1}\\
Undo the last commit, discard the changes &
  \code{git reset --hard HEAD\textasciitilde 1}\\
Undo a commit that is already pushed & \code{git revert <commit>}\\
Set aside unfinished work temporarily & \code{git stash} /
  \code{git stash pop}\\
\bottomrule
\end{tabularx}

\begin{gotcha}{\code{reset --hard} deletes work, permanently}
\code{git reset --hard} rewrites the working directory to match the
target commit and discards uncommitted changes with no confirmation.
\code{git revert} is the safe alternative once a commit has been shared
with anyone else: it adds a new commit that undoes the old one, instead
of rewriting history.
\end{gotcha}

\section{\file{.gitignore}}
\label{sec:commands:gitignore}

A plain text file, one pattern per line, listing what \Git{} should
never track: build products, logs, editor files, credentials.

\begin{txtcode}
# simulation and build output
*.log
*.vcd
work/
obj_dir/

# editor and OS clutter
.vscode/
.DS_Store
\end{txtcode}

Patterns are read from every \file{.gitignore} in the tree; a leading
\code{/} anchors a pattern to that folder, and a trailing \code{/}
matches directories only. A file already tracked before it was added to
\file{.gitignore} keeps being tracked; remove it once with
\code{git rm --cached <file>}.

\section{Checklist}
\label{sec:commands:checklist}

\begin{itemize}
  \item \code{status} and \code{diff} before every \code{add}; know
        what you are about to commit.
  \item Small, focused commits with an imperative-mood message beat one
        giant commit at the end of the day.
  \item \code{switch}/\code{restore} for everyday work; save
        \code{checkout} for detached-HEAD or older scripts.
  \item \code{revert} once a commit is shared; \code{reset --hard} only
        on commits nobody else has seen.
  \item Add a \file{.gitignore} before the first commit, not after.
\end{itemize}
